% Lösungen Einheit 1 von 4 – Erwartungswert berechnen und deuten: die Verteilung beschaffen (Tabelle, Sachtext, Baumpfade, Kosten je Ausgang; fehlende Wahrscheinlichkeit über die Summe 1), die gewichtete Summe bilden (bei Anzahlen n · p), das Ergebnis deuten (zusammenbau v0.1)
\einheitenkopf{Einheit 1 von 4 · Erwartungswert berechnen und deuten: die Verteilung beschaffen (Tabelle, Sachtext, Baumpfade, Kosten je Ausgang; fehlende Wahrscheinlichkeit über die Summe 1), die gewichtete Summe bilden (bei Anzahlen n · p), das Ergebnis deuten}

\erg{9}{a) Auszahlung \quad b) fehlend $1 - 0{,}5 - 0{,}3 = 0{,}2$; $E(X) = 0 \cdot 0{,}5 + 1 \cdot 0{,}3 + 2 \cdot 0{,}2 = 0{,}7$ \quad c) $E(X) = 50 \cdot 0{,}02 + 5 \cdot 0{,}1 + 0 = 1{,}50$ (Euro) \quad d) $E(\text{Auszahlung}) = 0{,}25 \cdot 10 = 2{,}50$ (Euro), weniger als der Einsatz; Gewinn des Anbieters $= 3 - 2{,}50 = 0{,}50$ (Euro) je Spiel \quad e) mittlerer Preis $= 1 + \frac{1}{2} \cdot 2 + \frac{1}{2} \cdot 4 + 2 = 6$ (Euro); Einnahmen $= 120 \cdot 6 = 720$ (Euro); $1\,000 : 6 \approx 166{,}7$, aufgerundet $n = 167$ Tüten \quad f) Der zweite Summand gehört zu genau einer roten Kugel: 4 ist die Auszahlung in Euro, der Faktor 2 zählt die zwei Reihenfolgen rot–blau und blau–rot, $0{,}3 \cdot 0{,}7$ ist die Wahrscheinlichkeit eines solchen Pfades}
\erg{10}{a) Erwartungswert $150 \cdot 0{,}06 = 9$; Abweichung $\frac{12 - 9}{9} \approx 33{,}3\,\%$ über dem Erwartungswert}
\erg{11}{a) $E(T) = 12 \cdot \frac{1}{3} = 4$; Augensumme $= 12 \cdot 1 + 3 \cdot T$, Erwartungswert $= 12 + 3 \cdot 4 = 24$}
\erg{12}{a) Annahme: in der letzten Gruppe genau 4 Personen; mittlere Belegung $= 1 \cdot 0{,}12 + 2 \cdot 0{,}46 + 3 \cdot 0{,}30 + 4 \cdot 0{,}12 = 2{,}42$; Zelte $= 2\,420 : 2{,}42 = 1\,000$}
\erg{13}{a) Fehler: Ben deutet die Auszahlung als Gewinn und rechnet den Einsatz nicht gegen; Gewinn $= 1{,}80 - 2 = -0{,}20$ (Euro), er verliert im Mittel 20 Cent je Spiel \quad b) Der Erwartungswert ist der Durchschnitt je Durchführung auf lange Sicht, ein Mittelwert der Werte, und muss deshalb selbst kein Wert sein, den die Zufallsgröße annehmen kann. \quad c) Säulen über 0, 1, 2, 3 mit den Höhen 0{,}1; 0{,}4; 0{,}3; 0{,}2 \quad d) mittlere Zahlung $= 800 \cdot 0{,}04 + 150 \cdot 0{,}1 = 47$ (Euro); Überschuss $= 60 - 47 = 13$ (Euro) je Vertrag}

13c)

\saeulen{0/0.1,1/0.4,2/0.3,3/0.2}{0.5}{0.1}{$P(X = x)$}

